% coord_R2_cart.tex
%
% Copyright (C) 2022--2026 José A. Navarro Ramón <janr.devel@gmail.com>
% Licencia del código GPLv2
% Licencia Creative Commons Recognition Non-Commercial Share-alike.
% (CC-BY-NC-SA)

\chapter{Sistema de coordenadas cartesianas ortogonales en
  \mathinhead{\symbb{R}^2}{espr2}}

Este capítulo tiene los siguientes objetivos:
\begin{enumerate}
\item Presentar el sistema de coordenadas ortogonales o rectantulares en el plano
  euclídeo $\symbb{R}^2$, como fundamento para el resto de sistemas de coordenadas que
  se estudiarán posteriormente.
\item Introducir los conceptos de base natural y de métrica en este sistema, conceptos
  que serán más importantes en otros sistemas de coordenadas.
\item Recordar la expresión del resultado de distintas operaciones de cálculo en este
  sistema.
\end{enumerate}

\section{Ideas básicas}
Vector de posición de cualquier punto del espacio $\symbb{R}^2$ en coordenadas cartesianas
\begin{equation}\label{eq:R2-r_rect}
  \vvv{r}
  = x\,\uvec{\i} + y\,\uvec{\j}
\end{equation}
\begin{figure}[ht]
  % Escala
  \def\scl{1}
  %
  \pgfmathsetmacro{\EJESEXTRA}{-0.2}
  \pgfmathsetmacro{\EJESLONG}{2.3}
  % Eje x
  \pgfmathsetmacro{\XMLONG}{\EJESEXTRA}
  \pgfmathsetmacro{\XPLONG}{\EJESLONG}
  % Eje y
  \pgfmathsetmacro{\YMLONG}{\EJESEXTRA}
  \pgfmathsetmacro{\YPLONG}{\EJESLONG}
  % Versores
  \pgfmathsetmacro{\VMOD}{0.6}
  % Vector P
  \pgfmathsetmacro{\PMOD}{2.5}
  \pgfmathsetmacro{\PANG}{40}
  % Fondo
  \pgfmathsetmacro{\HORZ}{0.25}
  \pgfmathsetmacro{\VERT}{0.25}
  % 
  \centering
  \begin{tikzpicture}[%
    scale=\scl,
    every node/.style={black,font=\small},
    eje/.style={->},
    proyeccion/.style={black!20, line width=0.4},
    vector/.style={-{Latex[round]}, shorten >=1.2pt, line width=.8pt,\colorV},
    versor/.style={-{Latex[round, width=4pt, length=6pt]}, line width=1pt, \colorVred},
    textooriginal/.style={\colorTorig},
    pcirculo/.style={fill=\colorPorig, draw=black},
    angulogirado/.style={fill=\colorAngGreen, draw=\colorGirodos},
    background/.style={%
      line width=\bgborderwidth,
      draw=\bgbordercolor,
      fill=\bgcolor,
    },
    ]
    % COORDENADAS
    % Origen
    \coordinate (O) at (0,0);
    % Extremo izdo del eje x e inferior del eje y
    \coordinate (xini) at (\XMLONG cm,0);
    \coordinate (yini) at (0,\YMLONG cm);
    \coordinate (ifin) at (\VMOD cm, 0);
    \coordinate (jfin) at (0, \VMOD cm);
    % Eje x
    \path[save path=\ejex] (xini) -- coordinate[pos=0.55] (xtexto) (\XPLONG cm, 0)
    coordinate (xfin);
    % Eje y
    \path[save path=\ejey] (yini) -- coordinate[pos=0.55] (ytexto) (0, \YPLONG cm)
    coordinate (yfin);
    % Versores i y j
    \path[save path=\i] (O) -- coordinate[pos=0.5] (itexto) (ifin);
    \path[save path=\j] (O) -- coordinate[pos=0.45] (jtexto) (jfin);    
    % Vector
    \path[save path=\vector] (O) -- coordinate[pos=0.5] (pmid) (\PANG:\PMOD cm)
    coordinate (P);
    % Proyección de P en los ejes
    \path[save path=\proyx] (P) |- (xfin);
    \path[save path=\proyy] (P) -| (yfin);

    % DIBUJOS
    % Proyección de P en los ejes
    \draw[proyeccion, use path=\proyx];
    \draw[proyeccion, use path=\proyy];
    % Ejes
    % x
    \draw[eje, use path=\ejex];
    \node[right] at (xfin) {$x$};
    \node[below] at (xtexto) {$x$};
    % y
    \draw[eje, use path=\ejey];
    \node[above] at (yfin) {$y$};
    \node[left] at (ytexto) {$y$};
    % Versores
    \draw[versor, use path=\i];
    \node[below, \colorTred] at (itexto) {\scriptsize $\uvec{\i}$};
    \draw[versor, use path=\j];
    \node[left, \colorTred] at (jtexto) {\scriptsize $\uvec{\j}$};
    % Vector y punto  P
    \draw[vector, use path=\vector];
    \node[above left=-3pt and -1pt] at (pmid) {$\vvv{r}$};
    \filldraw[fill=green, draw=black] (P) circle[radius=1.2pt];
    \node[above right=0pt and -5pt] at (P) {$P(x,y)$};
    
    % Origen
    \filldraw (O) circle [radius=.2pt];
    % Fondo amarillo
    \coordinate (SW) at ($(current bounding box.south west) + (-\HORZ cm,-\VERT cm)$);
    \coordinate (NE) at ($(current bounding box.north east) + (\HORZ cm,\VERT cm)$);    
    \begin{scope}[on background layer]
      \draw[background] (SW) rectangle (NE);
    \end{scope}
  \end{tikzpicture}
  \caption{Coordenadas cartesianas de un punto $P$ del plano (en verde) y vectores
    unitarios (en rojo). Se puede comprobar que el vector de posición del punto es
    $\vvv{r} = x\uvec{i} + y\uvec{j}$.}
  \label{fig:R2-coord_rect-ij}
\end{figure}


\subsection{Base natural}
\subsubsection{Base natural y vectores unitarios}
Si derivamos el vector de posición anterior \eqref{eq:R2-r_rect}, con respecto a
cada una de las coordenadas, obtenemos los vectores de una base, llamada \emph{natural},
en estas coordenadas. Estos vectores llevan información concreta acerca de cómo varia la
posición en el espacio cuando variamos cada una de las coordenadas ---tal y como
sugieren las derivadas---, lo que será importante para obtener la métrica en estas
coordenadas.

Como
\begin{equation}
  \vvv{r} = x\uvec{\i} + y\uvec{\j}
\end{equation}

Entonces, derivando se obtiene
\begin{subequations}
  \begin{align}\label{eq:R2-cart-ex}
    \xhat{e}_x
    &=
      \dfrac{\partial\vvv{r}}{\partial x}
      = \uvec{\i}
      \longrightarrow
    |\xhat{e}_x| = |\uvec{\i}| = 1\\
    \label{eq:R2-cart-ey}
    \xhat{e}_y
    &= \dfrac{\partial\vvv{r}}{\partial y}
      = \uvec{\j}
      \longrightarrow
    |\xhat{e}_y| = |\uvec{\j}| = 1
  \end{align}
\end{subequations}
estos vectores son unitarios aunque, en general, no tiene por qué ser así en otros
sistemas.

Como hemos dicho, estos vectores forman la denominada \emph{base natural} en cartesianas
que, en este sistema coincide con los vectores unitarios tradicionales $\uvec{\i}$ y
$\uvec{\j}$,
aunque en otros sistemas de coordenadas no tienen por qué ser unitarios
\begin{equation}\label{eq:R2-cart-base}
  B=\set{\xhat{e}_x, \xhat{e}_y} = \set{\uvec{\i},\, \uvec{\j}}
\end{equation}

La métrica en cada sistema de coordenadas define la estructura del espacio, en
particular, sin ella no se pueden medir ni distancias ni ángulos. Es como si
proporcionara un sistema de reglas y goniómetros con los cuales medir.

De la equivalencia de la base natural con la unitaria, se puede inferir que la métrica en
coordenadas cartesianas rectangulares en $\symbb{R}^2$ se representará como una matriz
unidad (en el apartado siguiente,\ref{ssect:R2-cart-metrica-factores_escala}, se
calculará con detalle).

\subsection{Métrica y factores de escala}
\label{ssect:R2-cart-metrica-factores_escala}

\subsubsection{Producto escalar de los vectores de la base natural}
Los vectores de la base fundamental son vectores unitarios.
La base normalizada en coordenadas cartesianas en este espacio está formada por los vectores $\uvec{\i}$ y $\uvec{\j}$, en los ejes $x$ e $y$, respectivamente.
A continuación resumimos el valor de los productos escalares entre ellos, tal y como
aprendimos en cálculo vectorial fundamental en coordenadas cartesianas o en física
elemental
\begin{subequations}
  \begin{align}
    \label{eq:R2-cart-ii_prod_escalar}
    \xhat{e}_x\cdot\xhat{e}_x = \uvec{\i}\cdot\uvec{\i}
    &= |\uvec{\i}|\,|\uvec{\i}|\cos 0 = 1\\
    \label{eq:R2-cart-ij_prod_escalar}    
    \xhat{e}_x\cdot\xhat{e}_y = \uvec{\i}\cdot\uvec{\j}
    &= |\uvec{\i}|\,|\uvec{\j}|\cos\pi/2 = 0\\
    \label{eq:R2-cart-ji_prod_escalar} 
    \xhat{e}_y\cdot\xhat{e}_x = \uvec{\j}\cdot\uvec{\i}
    &= |\uvec{\j}|\,|\uvec{\i}|\cos\pi/2 = 0\\
    \label{eq:R2-cart-jj_prod_escalar}
    \xhat{e}_y\cdot\xhat{e}_y = \uvec{\j}\cdot\uvec{\j}
    &= |\uvec{\j}|\,|\uvec{\j}|\cos 0 = 1
  \end{align}
\end{subequations}

\subsubsection{Métrica y factores de escala}
La métrica en cartesianas se calcula utilizando los vectores de la base natural,
calculada en \eqref{eq:R2-cart-ex} y \eqref{eq:R2-cart-ey}.
Los productos escalares se calcularon en \eqref{eq:R2-cart-ii_prod_escalar},
\eqref{eq:R2-cart-ij_prod_escalar}, \eqref{eq:R2-cart-ji_prod_escalar} y
\eqref{eq:R2-cart-jj_prod_escalar}
\begin{equation}\label{eq:R2-cart-metrica}
  \mmm{g}_{ij}
  = \begin{pmatrix}
    g_{11} & g_{12}\\
    g_{21} & g_{22}
    \end{pmatrix}
  = \begin{pmatrix}
      \xhat{e}_x\cdot\xhat{e}_x & \xhat{e}_x\cdot\xhat{e}_y\\
      \xhat{e}_y\cdot\xhat{e}_x & \xhat{e}_y\cdot\xhat{e}_y
  \end{pmatrix}
  = \begin{pmatrix}
    1 & 0\\
    0 & 1
  \end{pmatrix}
\end{equation}

La raíz cuadrada de los elementos diagonales de la métrica en un sistema de coordenadas ortogonales, como el que nos ocupa, se llaman \emph{factores de escala} o \emph{coeficientes métricos}. En coordenadas cartesianas rectangulares serían
\begin{align*}
  h_x &= \sqrt{g_{11}} = \sqrt{1} = 1\\
  h_y &= \sqrt{g_{22}} = \sqrt{1} = 1  
\end{align*}

Fijémonos en la métrica \eqref{eq:R2-cart-metrica}:
\begin{itemize}
\item Los elementos no diagonales son nulos. Esto se interpreta como que los vectores
  $\xhat{e}_x$ y $\xhat{e}_y$ son independientes, esto es, son ortogonales.
  Esto implica que un cambio en una coordenada no afecta a la otra. Así, si nos
  desplazamos una distancia cualquiera en el eje $x$, no habremos recorrido ninguna
  distancia en el eje $y$, y viceversa, por lo que $h_{xy}=h_{yx}=0$.
\item Los elementos diagonales valen la unidad. Cada elemento diagonal de la métrica que
  vale la unidad implica que un cambio en la coordenada respectiva implica una distancia
  recorrida de una unidad en el espacio. En realidad, el espacio recorrido es el factor
  de escala correspondiente (la raíz cuadrada del elemento diagonal correspondiente),
  pero en este caso da lo mismo
  \begin{subequations}
    \begin{align}
      \Delta s &= h_x \Delta x = \sqrt{1}\, \Delta x = \Delta x\\
      \Delta s &= h_y \Delta y = \sqrt{1}\, \Delta y = \Delta y
    \end{align}
  \end{subequations}
  Para adquirir algo de intuición, imaginemos que nos desplazamos cualquier distancia en
  el eje $x$, digamos un metro; es fácil entender que en el espacio nos habremos
  desplazado esa misma distancia. Por tanto es lógico razonar que el factor de escala
  $h_x$ vale 1. Lo mismo podemos decir para el factor de escala en el eje $y$, $h_y=1$.
\end{itemize}

\subsubsection{Revisión del producto escalar}
Dijimos que la métrica nos permite calcular distancias. En concreto, el cuadrado del
módulo de un vector $\vvv{A}$ se mide realizando el producto escalar del vector por sí
mismo.
El producto escalar se encarga de esto, y se puede ampliar mediante el cálculo
matricial, utilizando la métrica (en otros sistemas de coordenadas veremos que es más
importante, mientras que en cartesianas se puede obviar)
\[
  \vvv{v}\cdot\vvv{w}
  =
  \vvv{v}^\transp \mmm{g}_{ij} \vvv{w}
  =
  \begin{pmatrix}
    v_x &  v_y
  \end{pmatrix}
  \begin{pmatrix}
    1 & 0\\
    0 & 1
  \end{pmatrix}
  \begin{pmatrix}
    w_x\\
    w_y
  \end{pmatrix}
  =
  v_xw_x + v_y w_y
\]
  
Así, el cuadrado del módulo de un vector en cartesianas es
\[
  v^2
  =
  \|\vvv{v}\|^2
  =
  \vvv{v}\cdot\vvv{v}
  =
  \vvv{v}^\transp \mmm{g}_{ij} \vvv{v}
  =
  \begin{pmatrix}
  v_x &  v_y
\end{pmatrix}
\begin{pmatrix}
  1 & 0\\
  0 & 1
\end{pmatrix}
\begin{pmatrix}
  v_x\\
  v_y
\end{pmatrix}
=
v_x^2 + v_y^2
\]
que coincide con el valor que obteníamos en el cálculo vectorial básico.

\subsection{Gradiente, divergencia y laplaciana}
Presentamos un resumen de estas operaciones, tal y como se presentan en un curso básico
de cálculo
\begin{itemize}
\item El gradiente de un campo escalar $\phi(x,y)$ es un vector
\begin{equation}
  \vvv{\nabla}\phi
  = \dfrac{\partial \phi}{\partial x}\,\uvec{\i}
  + \dfrac{\partial \phi}{\partial y}\,\uvec{\j}
  = \begin{pmatrix}
    \partial \phi/\partial x\\
    \partial \phi/\partial y
    \end{pmatrix}
\end{equation}
El vector gradiente de $\phi$ evaluado en un punto genérico $(x,y)$ del dominio de $\phi$
indica la dirección en la cual el campo $\phi$ varía más rápidamente y su módulo
representa el ritmo de variación de $\phi$ en la dirección de dicho vector gradiente.

El operador gradiente se comporta como un vector bajo rotaciones, decimos esto porque
cuando opera frente a un campo escalar, se obtiene otro formado por vectores
\begin{equation}
  \vvv{\nabla}
  = \dfrac{\partial}{\partial x}\,\uvec{\i}
  + \dfrac{\partial}{\partial y}\,\uvec{\j}
  = \begin{pmatrix}
    \partial/\partial x\\
    \partial/\partial y
    \end{pmatrix}
\end{equation}

\item La divergencia de un campo vectorial $\vvv{A}=(A_x, A_y)$ es un escalar
  \begin{equation}\label{eq:R2-divergencia-cartesianas}
    \vvv{\nabla}\cdot\vvv{A}
    = \begin{pmatrix}
      \partial/\partial x & \partial/\partial y      
    \end{pmatrix}
    \begin{pmatrix}
      1 & 0\\
      0 & 1
    \end{pmatrix}
    \begin{pmatrix}
      A_x\\
      A_y
    \end{pmatrix}
    = \dfrac{\partial A_x}{\partial x}
    + \dfrac{\partial A_y}{\partial y}
  \end{equation}
  La divergencia mide la diferencia entre el flujo saliente y el flujo entrante de un
  campo vectorial sobre la superficie que rodea a un volumen de control, por tanto, si el
  campo tiene ''fuentes'' la divergencia será positiva, y si tiene "sumideros", la
  divergencia será negativa. La divergencia mide la rapidez neta con la que se conduce la
  materia al exterior de cada punto, y en el caso de ser la divergencia idénticamente
  igual a cero, describe al flujo incompresible del fluido, llamado también campo
  solenoidal.  
\item La laplaciana de un campo escalar $f(x,y)$ es la divergencia del gradiente y es un
  escalar
  \begin{equation}\label{eq:R2-gradiente-cartesianas}
    \Delta \phi
    =
    \nabla^2 \phi
    = \vvv{\nabla}\cdot\vvv{\nabla}\phi
    = \dfrac{\partial^2 \phi}{\partial x^2}
    + \dfrac{\partial^2 \phi}{\partial y^2}
  \end{equation}
Mide la concentración de un campo en un punto. Si es negativo, indica una tendencia a la
concentración (fuente); si es positivo, indica dispersión.

El operador laplaciano es
\begin{equation}
  \nabla^2
  = \dfrac{\partial^2}{\partial x^2}
  + \dfrac{\partial^2}{\partial y^2}
\end{equation}
\end{itemize}

\subsection{Elemento de área}
El elemento diferencial de área en cartesianas es
\begin{equation}
  d^2a = dx\,dy
\end{equation}
\begin{figure}[ht]
  % Escala
  \def\scl{1}
  %
  \pgfmathsetmacro{\EJESEXTRA}{-0.4}
  \pgfmathsetmacro{\EJESLONG}{2.3}
  % Eje x
  \pgfmathsetmacro{\XMLONG}{\EJESEXTRA}
  \pgfmathsetmacro{\XPLONG}{\EJESLONG}
  % Eje y
  \pgfmathsetmacro{\YMLONG}{\EJESEXTRA}
  \pgfmathsetmacro{\YPLONG}{\EJESLONG}
  % dx, dy
  \pgfmathsetmacro{\LONG}{0.9}
  \pgfmathsetmacro{\DX}{\LONG}
  \pgfmathsetmacro{\DY}{\LONG}
  % Posición del elemento de área
  \pgfmathsetmacro{\PX}{0.5}
  \pgfmathsetmacro{\PY}{1.0}
  % Fondo
  \pgfmathsetmacro{\HORZ}{0.25}
  \pgfmathsetmacro{\VERT}{0.25}
  % 
  \centering
  \begin{tikzpicture}[%
    scale=\scl,
    every node/.style={black,font=\small},
    eje/.style={->},
    proyeccion/.style={black!20, line width=0.4},
    area/.style={fill=green!40, draw=black!20},
    textooriginal/.style={\colorTorig},
    background/.style={%
      line width=\bgborderwidth,
      draw=\bgbordercolor,
      fill=\bgcolor,
    },
    ]
    % COORDENADAS
    % Origen
    \coordinate (O) at (0,0);
    % Extremo izdo del eje x e inferior del eje y
    \coordinate (xini) at (\XMLONG cm,0);
    \coordinate (yini) at (0,\YMLONG cm);
    % Eje x
    \path[save path=\ejex] (xini) -- (\XPLONG cm, 0) coordinate (xfin);
    % Eje y
    \path[save path=\ejey] (yini) -- (0, \YPLONG cm) coordinate (yfin);
    % Punto P
    \coordinate (P) at (\PX, \PY);
    %
    \path (P) -- coordinate[pos=0.5] (dxtext) ++(\DX, 0) coordinate (PDX);
    \path (P) -- coordinate[pos=0.5] (dytext) ++(0, \DY) coordinate (PDY);
%    % Vector
%    \path[save path=\vector] (O) -- coordinate[pos=0.5] (pmid) (\PANG:\PMOD cm)
%    coordinate (P);
%    % Proyección de elemento de área en los ejes
    \path[save path=\proyPx] (P) |- coordinate (Px) (xfin);
    \path[save path=\proyPy] (P) -| coordinate (Py) (yfin);
    \path[save path=\proyPhorx] (PDX) |- coordinate (Phorx) (xfin);
    \path[save path=\proyPverty] (PDY) -| coordinate (Pverty) (yfin);
    % Texto de las proyecciones en los ejes
    \path (Px) -- coordinate[pos=0.5] (dxtexto) (Phorx);
    \path (Py) -- coordinate[pos=0.5] (dytexto) (Pverty);
    
    % DIBUJOS
    % Proyección de elemento de área en los ejes
    \draw[proyeccion, use path=\proyPx];
    \draw[proyeccion, use path=\proyPy];
    \draw[proyeccion, use path=\proyPhorx];
    \draw[proyeccion, use path=\proyPverty];
    % Texto dx y dy en los ejes
    \node[below] at (dxtexto) {$dx$};
    \node[left] at (dytexto) {$dy$};
    
    % Ejes
    % x
    \draw[eje, use path=\ejex];
    \node[right] at (xfin) {$x$};
    %\node[below] at (xtexto) {$x$};
    % y
    \draw[eje, use path=\ejey];
    \node[above] at (yfin) {$y$};
    % \node[left] at (ytexto) {$y$};
    % Punto P
    % \fill (P) circle[radius=1.2pt];
    % Elemento de área
    \filldraw[area] (P) rectangle coordinate (centro) ++(\DX, \DY);
    \node at (centro) {\footnotesize $d^2a$};

    % Origen
    \filldraw (O) circle [radius=.2pt];

    % Fondo amarillo
    \coordinate (SW) at ($(current bounding box.south west) + (-\HORZ cm,-\VERT cm)$);
    \coordinate (NE) at ($(current bounding box.north east) + (\HORZ cm,\VERT cm)$);    
    \begin{scope}[on background layer]
      \draw[background] (SW) rectangle (NE);
    \end{scope}
  \end{tikzpicture}
  \caption{Elemento de área en coordenadas cartesianas, $d^2a = dx dy$.}
  \label{fig:R2-rect-elemento_area}
\end{figure}

Este elemento se relaciona trivialmente con el jacobiano de la
\emph{transformación identidad} $f(x,y) \longrightarrow f(x,y)$
\[
  d^2a = \dfrac{\partial(x,y)}{\partial(x,y)} dx dy
  = \begin{vmatrix}
    \partial x/\partial x & \partial x/\partial y\\
    \partial y/\partial x & \partial y/\partial y
  \end{vmatrix}
  dx dy
  = \begin{vmatrix}
    1 & 0\\
    0 & 1
  \end{vmatrix}
  dx dy
  = 1\,dx dy
  = dxdy
\]






%%% Local Variables:
%%% mode: latex
%%% TeX-engine: luatex
%%% TeX-master: "../retazosmatematicas.tex"
%%% End:
